\documentclass[compacto,10pt]{pubciencias-notas}

% -- Paquetes adicionales
\usepackage{comandosEPN}

% -- Datos del libro
\autor{Departamento de Formación Básica}
\materia{Álgebra Lineal}
\nota{Semana 05 Clase 01}
\fecha{Semestre 2025-A}

\logouno[1.1cm]{Logo_EPN.eps}
\logodos[1.7cm]{Logo_DFB.eps}

\definecolor{colordef}{cmyk}{0.81,0.62,0.00,0.22}

\begin{document}

\encabezado

\vspace*{-8mm}
\setcounter{section}{1}
\section{Sistemas de ecuaciones}
\setcounter{subsection}{3}
\subsection{Sistemas homogéneos}

\begin{teo}\label{Sem03Clas02CarcSisHom}
    Sea $A \in \Mat[\K]{n}{n}$, el sistema homogéneo 
    \[
        Ax = 0
    \]
    tiene una solución no trivial si y solo si $A$ es singular. 
\end{teo}

\begin{comment}
\begin{proof}\hspace{0pt} 
Supongamos que $A$ es singular, entonces $\rang(A)  = \rang(A|0) < n$, lo que implica el sistema $Ax = 0$ es consistente con infinitas soluciones, por tanto, el sistema tiene al menos una solución no trivial.

Para el recíproco razonemos por reducción al absurdo, esto es, supongamos que $Ax = 0$ tiene una solución no trivial y que $A$ es no singular, lo que implica que $A^{-1}$ existe, de lo cual
\[
    Ax = 0 \Longleftrightarrow A^{- 1}(Ax) = A^{- 1}0 \Longleftrightarrow x = 0
\]
lo que contradice lo supuesto.
\end{proof}
\end{comment}

\begin{teo}\label{Sem03Clas02InvCat}
    Sea $A \in \Mat[\K]{n}{n}$. Se tiene que $A$ es no singular si y solo si el sistema lineal $Ax=b$ tiene una solución única para cada $b \in \Mat[\K]{n}{1}$.
\end{teo}

\begin{ejem}\hspace{0pt}\label{Sem03Clas02EjemPartB} 
Dado el sistema de ecuaciones lineales
\[
    \systeme[x_1 x_2]{2x_1 +x_2 =b_{1},
        x_1 +x_2 = b_{2}}
\]
que tiene solución única para todo $\begin{pmatrix}b_{1}\\b_{2}\end{pmatrix}\in\Mat{2}{1}$, se tiene que 
$A = \begin{pmatrix}2&1\\1&1\end{pmatrix}$ es no singular. En efecto, notemos que 
\[
    \begin{pmatrix}
    2&1&|&b_{1}\\
    1&1&|&b_{2}
    \end{pmatrix}
    \sim
    \begin{pmatrix}
    1&1&|&b_{2}\\
    0&-1&|&b_{1} - 2b_{2}
    \end{pmatrix}
\]
luego, $\rang(A) = \rang(A|b) = 2$, y, por tanto, el sistema $Ax = b$ posee solución única, que junto con el teorema \ref{Sem03Clas02InvCat} implica que $A$ es no singular.
\end{ejem}

\begin{ejer}\hspace{0pt} 
Considere al sistema de ecuaciones lineales
\[
    \systeme[x_1 x_2 x_3]{x_1 +2x_2 =1,
        x_2 +2x_3 = 1,
        x_3 = 1}
\]
Utilizando la inversa de la matriz de coeficientes muestre que el sistema posee solución única y encuéntrela.
\end{ejer}

\begin{prop}\label{Sem03Clas02EquivInv}
    Sea $A \in \Mat[\K]{n}{n}$, se tienen que las siguientes proposiciones son equivalentes:
    \begin{enumerate}
    \item[$a)$] 
        $A$ es no singular;
    \item[$b)$] 
        el sistema $Ax=0$ solamente tiene la solución trivial;
    \item[$c)$] 
        $A$ es equivalente por filas a $I_n$;
    \item[$d)$] 
        $\rang(A)=n$; y 
    \item[$e)$] 
        el sistema lineal $Ax=b$ tiene una solución única para cada $b \in \Mat[\K]{n}{1}$. 
    \end{enumerate}
\end{prop}

\begin{comment}
\begin{ejem} Dada la matriz 
\[
    A = \begin{pmatrix}
    2&1\\
    1&1
    \end{pmatrix}
\]
La matriz $A$ verifica a la proposición \ref{Sem03Clas02EquivInv}.

\begin{proof}\hspace{0pt} Considerando al ejemplo \ref{Sem03Clas02EjemPartA} notemos que los literales $a)$, $c)$ y $d)$ se verifican. Considerando al ejemplo \ref{Sem03Clas02EjemPartB} notemos que se cumple el literal $e)$. Por el literal $a)$ junto con el teorema \ref{Sem03Clas02CarcSisHom}, el literal $b)$ se verifica. Por lo tanto, la matriz $A$ verifica a la proposición \ref{Sem03Clas02EquivInv}.\qedhere
\end{proof}
\end{ejem}

\begin{ejer}\hspace{0pt} Construya un ejemplo de un sistema de ecuaciones lineales $3 \times 3$ que tenga solución, y compruebe que su matriz de coeficientes es singular.
\end{ejer}
\end{comment}

\begin{ejer}\hspace{0pt} 
Dada la matriz $A = \begin{pmatrix}a&b\\c&d\end{pmatrix}$, verificar que su inversa es  
\[
    A^{-1} = \dfrac{1}{ad - bc}\begin{pmatrix}d&-b\\-c&a\end{pmatrix} \text{   siempre que   } ad - bc \neq 0.
\]
\end{ejer}

\setcounter{section}{2}
\section{Determinantes}

En esta sección tomaremos $n\in \N$, con $n>0$, e $I = \{ 1, 2, \ldots, n\}$.

\subsection{Definiciones fundamentales}

\begin{defi}[Menor]
    Sea $A\in \Mat[\K]{n}{n}$ e $i,j\in I$. A la matriz de $\Mat[\K]{(n-1)}{(n-1)}$ que se obtiene luego de eliminar la fila $i$ y la columna $j$ de $A$ se la llama el menor $ij$ de $A$, y se denota por $A_{ij}$.
\end{defi}

\begin{advertencia}
    En la literatura se puede encontrar la notación de $M_{ij}$ para el menor de $ij$ de $A$.
\end{advertencia}

\begin{comment}
\begin{defi}[Menor principal]
    Sean $A\in \Mat[\K]{n}{n}$ y $k\in I$. A la matriz de $\Mat[\K]{k}{k}$ que se obtiene eliminar las $n-k$ últimas filas y columnas de $A$, se la llama el menor principal $k$ de $A$, denotado por $A_{k}$.
\end{defi}

\begin{advertencia}
    En la literatura se puede encontrar la notación de $M_{k}$ para el menor principal $k$ de $A$.
\end{advertencia}
\end{comment}

\begin{ejem}\hspace{0pt}

    Sea
    $A = \begin{pmatrix}
        -1 & 4 & 0 \\
        2 & 1 & -2 \\
        0 & 3 & 5
    \end{pmatrix}$, se tiene que:
    
    \[A_{12} = \begin{pmatrix}
        2 & -2 \\
        0 & 5
    \end{pmatrix}, \quad 
    A_{23} = \begin{pmatrix}
        -1 & 4 \\
        0 & 3
    \end{pmatrix} \quad \text{y} \quad 
    A_{31} = \begin{pmatrix}
        4 & 0 \\
        1 & -2
    \end{pmatrix}
    \]
    
    %\noindent Determine los menores $A_{32}$, $A_{13}$ y $A_{22}$.

\end{ejem}

\begin{ejer}
\item Sea
$B = \begin{pmatrix}
    1 & -3 & -1 \\
    0 & 1 & -5 \\
    4 & 2 & 0
\end{pmatrix}$. Determine los menores $B_{11}$, $B_{32}$ y $B_{21}$.
\end{ejer}


\begin{defi}[Determinante] 
    Sea $A\in \Mat[\K]{n}{n}$ se define el determinante de $A$, denotado por $\det(A)$ (o por $|A|$), de manera inductiva, como sigue:
    \begin{itemize}
    \item 
        Si $n=1$ y $A=(a_{11})$, entonces $\det(A) = a_{11}$.
    \item
        Si $n>1$, entonces
        \begin{align*}
            \det(A) 
            & = \sum_{k=1}^n a_{1k}(-1)^{1+k}\det(A_{1k})\\
            & = a_{11}\det(A_{11}) - a_{12}\det(A_{12}) 
            + \ldots +  (-1)^{1+n} a_{1n}\det(A_{1n}).
        \end{align*}
    \end{itemize}
\end{defi}

\begin{ejem}\hspace{0pt} \label{determinante3por3}

\begin{itemize}
\item 
    Sea $A$ una matriz de orden $2 \times 2$ de la forma
    \[ 
        A=
        \begin{pmatrix}
         a_{11} & a_{12}\\
         a_{21} & a_{22}
        \end{pmatrix},
    \]
    se tiene que
    \[
        A_{11} = (a_{22})
        \texty
        A_{12} = (a_{21}),
    \]
    por lo tanto
    \[
        \det(A_{11}) = a_{22}
        \texty
        \det(A_{12}) = a_{21},
    \]
    de esta forma,
    \begin{align*}
        \det(A) &= \sum_{k=1}^{2} a_{1k}(-1)^{1+k}\det(A_{1k})\\
        &= (-1)^{1+1}a_{11}\det(A_{11}) + (-1)^{1+2}a_{12}\det(A_{12})\\
        & = a_{11}a_{22} - a_{12}a_{21}
    \end{align*}
    es decir,
    \[
        \det(A) = \begin{vmatrix}
        a_{11} & a_{12} \\
        a_{21} & a_{22}
        \end{vmatrix}
        = a_{11}a_{22} - a_{12}a_{21}.
    \]
    
\item Sea $B = \begin{pmatrix}
        1 & 5 \\
        2 & 6
    \end{pmatrix}$, se tiene que $\det(B) = 6-10 = -4$

%\item Sea $C$ una matriz de orden $3 \times 3$ de la forma
 %        \[ C=
  %      \begin{pmatrix}
   %      c_{11} & c_{12} & c_{13}\\
    %     c_{21} & c_{22} & c_{23}\\
     %    c_{31} & c_{32} & c_{33}\\
      %  \end{pmatrix}
       % \]
        %\noindent el determinante de la matriz $C$ está dado por:
        %\begin{align*}
        %\det(C) &= \sum_{k=1}^{3} (-1)^{1+k}c_{1k}\det(C_{1k}) \\
        %&= (-1)^{1+1}c_{11}\det(C_{11}) + 
        %(-1)^{1+2}c_{12}\det(C_{12}) + 
        %(-1)^{1+3}c_{13}\det(C_{13})\\
        %&= 
        %c_{11}
        %\begin{vmatrix}
        %c_{22} & c_{23} \\
        %c_{32} & c_{33}
        %\end{vmatrix} - c_{12}
      %  \begin{vmatrix}
%        c_{21} & c_{23} \\
 %       c_{31} & c_{33} \\
  %      \end{vmatrix}+ c_{13}
   %     \begin{vmatrix}
    %    c_{21} & c_{22} \\
     %   c_{31} & c_{32}
      %  \end{vmatrix}.
    %\end{align*}
%\item
 %   Sea $D = \begin{pmatrix}
  %      -1 & 1 & 0 \\
   %     0 & 1 & -2 \\
    %    2 & 2 & 1
    %\end{pmatrix}$, se tiene %que:
    %\[
     %   \det(D) = -1 
      %  \begin{vmatrix}
       %     1 & -2 \\
        %    2 & 1
        %\end{vmatrix} -1
        %\begin{vmatrix}
         %   0 & -2 \\
          %  2 & 1
        %\end{vmatrix} +0
        %\begin{vmatrix}
         %   0 & 1 \\
          %  2 & 2
        %\end{vmatrix} = -1(5) -1(4) = -9.
    %\]
\end{itemize}
\end{ejem}

%\begin{ejer}\hspace{0pt} Calcular el determinante de las siguientes matrices:
%\begin{enumerate}
%\item[$a)$] $A = \begin{pmatrix}-4 & 2 \\-1 & 3\end{pmatrix}$
%\item[$b)$] $B = \begin{pmatrix}1&-3&-1\\-2&1&-2\\0&5&0\end{pmatrix}$
%\end{enumerate}
%\end{ejer}

\begin{defi}[Cofactor]
    Sea $A\in \Mat[\K]{n}{n}$ e $i,j\in I$. El cofactor $ij$ de $A$ denotado $C_{ij}$, está dado por
    \[
        C_{ij} = (-1)^{i+j} \det (A_{ij})
    \]
     donde $A_{ij}$ es el menor $ij$ de $A$.
\end{defi}

\begin{ejem} Dada la matriz $A = 
\begin{pmatrix}
-2 & -1 & -3 \\
0 & 1 & -1 \\
3 & 3 & 2
\end{pmatrix}$, se tiene que:
\begin{itemize}
\item El cofactor $32$ de A es:
\[C_{32} = (-1)^{5}
\begin{vmatrix}
-2 & -3 \\
0 & -1
\end{vmatrix} = -1(2)=-2\]
\item El cofactor $13$ de A es:
\[C_{13} = (-1)^{4}
\begin{vmatrix}
0 & 1 \\
3 & 3
\end{vmatrix} = 1(-3) = -3\]
\end{itemize}
\end{ejem}

\begin{ejer}\hspace{0pt} Dada la matriz 
\[
    A = 
    \begin{pmatrix}
        -1 & 1 & 0 \\
        0 & 2 & -1 \\
        3 & 2 & 5
    \end{pmatrix}.
\]
Determinar los cofactores $C_{23}$, $C_{13}$, $C_{21}$ y $C_{32}$.
\end{ejer}

%\begin{advertencia}
%    En la literatura, también se suele llamar menor al determinante de $A_{ij}$ en 
%    lugar de a la matriz, como lo haremos en este texto. Además, cofactor, se lo %suele denotar por $A_{ij}$.
%\end{advertencia}

\begin{teo}[Generalización cálculo del determinante]
    Sea $A \in \Mat[\K]{n}{n}$. Se tiene que para todo $i\in I$,
    \[
        \det(A) 
        = \sum_{k=1}^n a_{ik}C_{ik} 
        = a_{i1}C_{i1} + a_{i2}C_{i2} + \ldots + a_{in}C_{in}
    \]
    y
    \[
        \det(A) 
        = \sum_{k=1}^n a_{ki}C_{ki} 
        = a_{1i}C_{1i} + a_{2i}C_{2i} + \ldots + a_{ni}C_{ni}.
    \]
\end{teo}

\begin{advertencia}
    El lado derecho de las igualdades toma el nombre de expansión por cofactores del determinante de $A$, con respecto a la fila $i$ y con respecto a la columna $i$.
\end{advertencia}

\begin{ejem}\hspace{0pt}\label{Det_Ej1}
    Sea $A =
    \begin{pmatrix}
        -2 & -1 & -3 \\
        0 & 1 & -1 \\
        3 & 3 & 2
    \end{pmatrix}$.
    \noindent Calculemos el determinante de $A$ mediante la expansión por cofactores de $A$, con respecto a la fila $1$. Se tiene
    \[
        \det(A) = -2 
        \begin{vmatrix}
            1 & -1 \\
            3 & 2
        \end{vmatrix} +1
        \begin{vmatrix}
            0 & -1 \\
            3 & 2
        \end{vmatrix} -3
        \begin{vmatrix}
            0 & 1 \\
            3 & 3
        \end{vmatrix} = 2
    \]
    
    ahora, calculemos el determinante de $A$ mediante la expansión por cofactores de $A$, con respecto a la columna 2:
    \[
        \det(A) = 1 
        \begin{vmatrix}
            0 & -1 \\
            3 & 2
        \end{vmatrix} +1
        \begin{vmatrix}
            -2 & -3 \\
            3 & 2
        \end{vmatrix} -3
        \begin{vmatrix}
            -2 & -3 \\
            0 & -1
        \end{vmatrix} = 2
    \]
    
    \end{ejem}
    
    \begin{ejer}\hspace{0pt} Considere a la matriz del ejemplo \ref{Det_Ej1}, calcule su determinante mediante la expansión por cofactores de $A$, con respecto a la columna $3$.
\end{ejer}
\end{document}